\documentclass[10pt,a4paper]{article}

\usepackage[a4paper,margin=0.62in]{geometry}
\usepackage[T2A]{fontenc}
\usepackage[utf8]{inputenc}
\usepackage[russian]{babel}

\usepackage{amsmath,amssymb}
\usepackage{booktabs}
\usepackage{array}
\usepackage{graphicx}
\usepackage{xcolor}
\usepackage[most]{tcolorbox}
\usepackage{enumitem}
\usepackage{titlesec}
\usepackage{helvet}
\usepackage{longtable}
\usepackage{tabularx}
\usepackage{fancyhdr}
\usepackage{lastpage}
\usepackage{tikz}
\usepackage{pifont}
\usepackage{etoolbox}
\usepackage{multicol}
\usepackage{pagecolor}
\usepackage{colortbl}
\usepackage{ragged2e}

\renewcommand{\familydefault}{\sfdefault}

% ---------- COLORS ----------
\definecolor{mainblue}{HTML}{5C6F82}      % спокойный slate для основного документа
\definecolor{lightblue}{HTML}{EEF3F6}     % мягкий голубовато-серый
\definecolor{softgray}{HTML}{F7F4F2}      % теплый фон

\definecolor{brandcoral}{HTML}{F29B8F}    % цвет логотипа
\definecolor{brandcoraldeep}{HTML}{E5867B}% чуть глубже
\definecolor{brandblush}{HTML}{FDF0EC}    % очень светлый коралловый
\definecolor{brandrose}{HTML}{F8D8D0}     % промежуточный розовато-коралловый

\definecolor{goodgreen}{HTML}{4E8B62}
\definecolor{lightgreen}{HTML}{ECF7F0}
\definecolor{warnorange}{HTML}{D98A3A}
\definecolor{lightorange}{HTML}{FFF2E5}
\definecolor{badred}{HTML}{C65A54}
\definecolor{lightred}{HTML}{FDEDEC}

\definecolor{darktext}{HTML}{25313D}
\definecolor{mutedtext}{HTML}{6B7280}
\definecolor{linegray}{HTML}{DED8D4}

\definecolor{hero}{HTML}{F29B8F}
\definecolor{heroDark}{HTML}{E5867B}

% ---------- GENERAL ----------
\setlength{\parindent}{0pt}
\setlength{\parskip}{5pt}

\setlist[itemize]{
    leftmargin=1.35em,
    itemsep=3pt,
    topsep=4pt
}
\setlist[enumerate]{
    leftmargin=1.45em,
    itemsep=3pt,
    topsep=4pt
}

\color{darktext}

% ---------- SECTION STYLE ----------
\titleformat{\section}
  {\Large\bfseries\color{mainblue}}
  {\thesection.}{0.45em}{}

\titleformat{\subsection}
  {\normalsize\bfseries\color{mainblue}}
  {\thesubsection}{0.45em}{}

\titlespacing*{\section}{0pt}{14pt}{8pt}
\titlespacing*{\subsection}{0pt}{10pt}{5pt}

% ---------- HEADER / FOOTER ----------
\pagestyle{fancy}
\fancyhf{}
\setlength{\headheight}{18pt}
\fancyhead[L]{\textcolor{brandcoraldeep}{\textbf{Preggo}}}
\fancyhead[R]{\textcolor{mutedtext}{Business Cup Stage 2}}
\fancyfoot[R]{\textcolor{mutedtext}{\thepage/\pageref{LastPage}}}
\renewcommand{\headrulewidth}{0.5pt}
\renewcommand{\footrulewidth}{0pt}
\renewcommand{\headrule}{\hbox to\headwidth{\color{linegray}\leaders\hrule height \headrulewidth\hfill}}

% ---------- SMALL HELPERS ----------
\newcommand{\cmark}{\textcolor{goodgreen}{\ding{51}}}
\newcommand{\warnmark}{\textcolor{warnorange}{\ding{72}}}
\newcommand{\xmark}{\textcolor{badred}{\ding{55}}}

\newcommand{\pill}[2]{%
\begin{tcolorbox}[
    enhanced,
    boxrule=0pt,
    arc=10pt,
    colback=#1,
    left=7pt,right=7pt,top=3pt,bottom=3pt,
    width=\fittoheight
]
\centering \small\bfseries #2
\end{tcolorbox}
}

\newcommand{\bluequote}[1]{%
\begin{tcolorbox}[
    enhanced,
    breakable,
    colback=brandblush,
    colframe=brandcoraldeep,
    boxrule=0.85pt,
    arc=3pt,
    left=8pt,right=8pt,top=6pt,bottom=6pt
]
#1
\end{tcolorbox}
}

% ---------- CUSTOM BOXES ----------
\newtcolorbox{summarybox}[1][]{
    enhanced,
    breakable,
    colback=lightblue,
    colframe=mainblue,
    boxrule=1pt,
    arc=4pt,
    left=10pt,right=10pt,top=8pt,bottom=8pt,
    fonttitle=\bfseries,
    title=#1
}

\newtcolorbox{softbox}[1][]{
    enhanced,
    breakable,
    colback=brandblush,
    colframe=brandcoraldeep,
    boxrule=0.85pt,
    arc=4pt,
    left=9pt,right=9pt,top=7pt,bottom=7pt,
    fonttitle=\bfseries,
    title=#1
}

\newtcolorbox{greenbox}[1][]{
    enhanced,
    breakable,
    colback=lightgreen,
    colframe=goodgreen,
    boxrule=0.8pt,
    arc=4pt,
    left=9pt,right=9pt,top=7pt,bottom=7pt,
    fonttitle=\bfseries,
    title=#1
}

\newtcolorbox{orangebox}[1][]{
    enhanced,
    breakable,
    colback=lightorange,
    colframe=warnorange,
    boxrule=0.8pt,
    arc=4pt,
    left=9pt,right=9pt,top=7pt,bottom=7pt,
    fonttitle=\bfseries,
    title=#1
}

\newtcolorbox{redbox}[1][]{
    enhanced,
    breakable,
    colback=lightred,
    colframe=badred,
    boxrule=0.8pt,
    arc=4pt,
    left=9pt,right=9pt,top=7pt,bottom=7pt,
    fonttitle=\bfseries,
    title=#1
}

\newtcolorbox{metriccard}[3]{
    enhanced,
    width=\linewidth,
    boxrule=0pt,
    arc=5pt,
    colback=#1,
    left=9pt,right=9pt,top=8pt,bottom=8pt
}

% ---------- ITEM STYLE ----------
\renewcommand\labelitemi{\textcolor{mainblue}{\ding{228}}}
\renewcommand\labelitemii{\textcolor{mainblue}{\ding{118}}}

% ---------- HERO ----------
\newcommand{\heroblock}{
\begin{tcolorbox}[
    enhanced,
    colback=hero,
    colframe=hero,
    arc=6pt,
    boxrule=0pt,
    left=14pt,right=14pt,top=14pt,bottom=14pt,
    overlay={
        \begin{scope}
            \clip (frame.south west) rectangle (frame.north east);
            \fill[white,opacity=0.08] 
                ([xshift=-1cm,yshift=2.8cm]frame.south west)
                rectangle
                ([xshift=18cm,yshift=4.2cm]frame.south west);
            \fill[heroDark,opacity=0.18]
                ([xshift=8.5cm,yshift=-0.2cm]frame.south west)
                circle (4.2cm);
            \fill[heroDark,opacity=0.15]
                ([xshift=13.8cm,yshift=2.9cm]frame.south west)
                circle (3.1cm);
        \end{scope}
    }
]
\begin{minipage}[c]{0.18\textwidth}
    \centering
    \IfFileExists{Preggo.png}{
        \begin{tcolorbox}[
            colback=white,
            colframe=white,
            arc=7pt,
            boxrule=0pt,
            width=2.8cm,
            height=2.8cm,
            halign=center,
            valign=center,
            left=2pt,right=2pt,top=2pt,bottom=2pt
        ]
        \includegraphics[width=0.92\linewidth]{Preggo.png}
        \end{tcolorbox}
    }{
        \begin{tcolorbox}[
            colback=white,
            colframe=white,
            arc=7pt,
            boxrule=0pt,
            width=2.8cm,
            height=2.8cm,
            halign=center,
            valign=center
        ]
        \textcolor{brandcoraldeep}{\Large\textbf{P}}
        \end{tcolorbox}
    }
\end{minipage}
\hfill
\begin{minipage}[c]{0.77\textwidth}
    {\color{white}\fontsize{24}{27}\selectfont \textbf{Preggo Unit Economics}}\\[5pt]
    {\color{white!95}\large B2C FemTech / AI Pregnancy Health Subscription}\\[7pt]
    {\color{white!88}\normalsize Not just a tracker. A safety system for pregnancy.}
\end{minipage}
\end{tcolorbox}
}

% ---------- METRIC ROW ----------
\newcommand{\metricrow}{
\noindent
\begin{minipage}[t]{0.24\textwidth}
\begin{tcolorbox}[enhanced,colback=lightgreen,colframe=goodgreen,boxrule=0.8pt,arc=4pt,left=7pt,right=7pt,top=7pt,bottom=7pt]
{\small\textcolor{mutedtext}{LTV}}\\[3pt]
{\Large\bfseries \$114.42}
\end{tcolorbox}
\end{minipage}
\hfill
\begin{minipage}[t]{0.24\textwidth}
\begin{tcolorbox}[enhanced,colback=lightblue,colframe=mainblue,boxrule=0.8pt,arc=4pt,left=7pt,right=7pt,top=7pt,bottom=7pt]
{\small\textcolor{mutedtext}{CAC}}\\[3pt]
{\Large\bfseries \$34.83}
\end{tcolorbox}
\end{minipage}
\hfill
\begin{minipage}[t]{0.24\textwidth}
\begin{tcolorbox}[enhanced,colback=lightorange,colframe=warnorange,boxrule=0.8pt,arc=4pt,left=7pt,right=7pt,top=7pt,bottom=7pt]
{\small\textcolor{mutedtext}{LTV/CAC}}\\[3pt]
{\Large\bfseries 3.28x}
\end{tcolorbox}
\end{minipage}
\hfill
\begin{minipage}[t]{0.24\textwidth}
\begin{tcolorbox}[enhanced,colback=softgray,colframe=darktext,boxrule=0.8pt,arc=4pt,left=7pt,right=7pt,top=7pt,bottom=7pt]
{\small\textcolor{mutedtext}{Payback}}\\[3pt]
{\Large\bfseries 2.13 мес}
\end{tcolorbox}
\end{minipage}
}

% ---------- DOCUMENT ----------
\begin{document}

\heroblock

\vspace{6pt}
\metricrow

\vspace{8pt}

\begin{summarybox}[\textbf{Суть модели в одном абзаце}]
Мы строим Preggo как подписочный B2C digital health продукт для беременных женщин. Наша задача в этой модели --- не просто показать красивые цифры, а доказать, что экономика сходится: пользователя можно привлекать, монетизировать и масштабировать без того, чтобы бизнес сгорал на платном трафике.

\vspace{4pt}
\[
\boxed{\textbf{LTV = \$114.42}}
\qquad
\boxed{\textbf{CAC = \$34.83}}
\qquad
\boxed{\textbf{LTV/CAC = 3.28x}}
\qquad
\boxed{\textbf{Payback = 2.13 мес}}
\]

Это значит, что при нашей текущей логике продукт можно масштабировать, если мы удерживаем сильный annual-first pricing и не зависим только от paid acquisition.
\end{summarybox}

\section{Что именно мы моделируем}

Preggo --- не просто pregnancy tracker и не просто контент-приложение. Мы продаем более дорогую и более сильную ценность:

\begin{itemize}
    \item continuous monitoring во время беременности;
    \item AI-интерпретацию симптомов и сигналов риска;
    \item lab/document scanning;
    \item doctor-facing reports;
    \item alerts и понятную картину состояния.
\end{itemize}

\bluequote{\textbf{Именно поэтому наша монетизация строится не вокруг дешевого entertainment-продукта, а вокруг подписки на спокойствие, контроль и безопасность в критически важный период жизни.}}

\newpage

\section{Портрет пользователя и его поведение}

Чтобы наши assumptions не выглядели абстрактными, мы фиксируем конкретный портрет пользователя, под которого строится вся модель.

\begin{softbox}[\textbf{Пользовательский портрет}]
\textbf{Emily, 30, first pregnancy, USA}

\begin{itemize}
    \item первая беременность, 8--16 неделя;
    \item уже тратит заметные деньги на врачей, анализы и обследования;
    \item тревожится из-за того, что между визитами к врачу у нее нет контроля;
    \item постоянно ищет информацию, но не доверяет разрозненным источникам;
    \item ей не нужны сравнения вроде ``ребенок размером с авокадо'' --- ей нужно понимать свое состояние и риски;
    \item здоровье ребенка для нее абсолютный приоритет, поэтому она готова платить за инструмент, который дает спокойствие, контроль и раннее замечание рисков.
\end{itemize}
\end{softbox}

\textbf{Связь с продуктом и pricing:}  
Мы сознательно не конкурируем с дешевыми pregnancy apps, которые продают контент и вовлечение. Мы строим продукт для пользователей вроде Emily, которым недостаточно ``интересных фактов'' --- им нужно понимать свое состояние и заранее замечать риски. Именно поэтому наша модель предполагает более высокий ARPU и premium pricing.

\textbf{Важно:}  
Все ключевые assumptions в нашей модели --- conversion, pricing, retention и каналы привлечения --- строятся именно вокруг такого пользователя. Это не массовый пользователь всех pregnancy apps, а более узкий сегмент с высокой готовностью платить за контроль, безопасность и предсказуемость.

\section{Откуда берутся assumptions}

Все исходные числа вынесены в Excel в лист \textbf{Assumptions}. Это важно: мы специально отделили assumptions от итоговых метрик, чтобы не было ``магических чисел''.

Ниже логика по каждому ключевому input'у.

\subsection{Цены: \$159 annual и \$15 monthly}

\begin{softbox}[\textbf{Базовые цены продукта}]
\begin{itemize}
    \item \textbf{Annual plan = \$159}
    \item \textbf{Monthly plan = \$15}
\end{itemize}
\end{softbox}

Почему такие расценки:
\begin{itemize}
    \item беременность --- это не одноразовая проблема на 1 неделю;
    \item ценность продукта высокая: здоровье матери + безопасность ребенка;
    \item annual-план дает лучший cash flow, лучше удерживает пользователя и лучше бьется с жизненным циклом беременности.
\end{itemize}

В Excel это лежит в листе \textbf{Assumptions} как основные editable inputs.

\subsection{Почему annual share = 80\%, а monthly share = 20\%}

В этот \textbf{assumption} мы закладываем, что большинство платящих пользователей выберут annual, потому что:

\begin{itemize}
    \item use case долгий: беременность тянется месяцами;
    \item core value у нас не развлекательная, а тревожно-медицинская;
    \item в такой категории многие готовы платить upfront, если получают ощущение контроля и safety;
    \item monthly-план нужен как entry option для менее уверенных пользователей, но мы не строим экономику на нем.
\end{itemize}

То есть 80/20 --- это \textbf{наша рабочая гипотеза для base-case}, а не абсолютная истина. Именно поэтому в Excel есть лист \textbf{Sensitivity}, где annual mix отдельно стресс-тестится.

\[
\text{Annual Share} = 80\%, \qquad \text{Monthly Share} = 20\%
\]

Monthly share в Excel считается как:
\[
\text{Monthly Share} = 1 - \text{Annual Share}
\]

\subsection{Почему monthly users платят 6 месяцев}

Тоже assumption. Мы не говорим, что любой monthly user сидит с нами весь срок беременности. Мы берем \textbf{умеренно-консервативную оценку}:

\begin{itemize}
    \item кто-то купит monthly ``попробовать'';
    \item кто-то зайдет позже, не в первом месяце беременности;
    \item кто-то может отвалиться раньше;
    \item поэтому в base-case мы берем не 9 месяцев и не 12, а \textbf{6 месяцев}.
\end{itemize}

Это не агрессивная цифра, а скорее осторожная.  
В Excel это input \textbf{Months on monthly plan = 6} на листе \textbf{Assumptions}.

\subsection{Почему platform fee = 15\%}

Мы закладываем комиссию платформы и платежной инфраструктуры на уровне:
\[
15\%
\]

Это агрегированное допущение под mobile subscription продукт: app store fee + processing + небольшой operational friction. Мы сознательно не ставим слишком низкое число, чтобы не сделать модель искусственно красивой.

\subsection{Почему variable cost per paid user = \$9}

Это переменные расходы на обслуживание одного платящего пользователя:

\begin{itemize}
    \item AI inference;
    \item cloud/storage;
    \item document processing;
    \item basic support;
    \item backend usage tied to actual paying usage.
\end{itemize}

То есть это не fixed cost команды и не burn компании, а именно \textbf{variable delivery cost per paid user}. Важно разделять эти вещи, иначе LTV получится грязным.

\subsection{Почему CPI = \$3.8}

\textbf{CPI} --- стоимость одной установки приложения с платного трафика.

\begin{softbox}[\textbf{Base-case}]
\[
CPI = \$3.8
\]
\end{softbox}

Пояснение:
\begin{itemize}
    \item health / pregnancy --- это не самый дешевый consumer category;
    \item аудитория узкая и качественная;
    \item продукт дорогой по perceived value, значит можно терпеть не самый низкий CPI;
    \item при этом \$3.8 --- это не запредельно высокий CPI для эмоциональной мобильной ниши.
\end{itemize}

Это base-case для paid social/performance логики. В \textbf{Sensitivity} мы отдельно проверяем, что будет при худшем CPI.

\subsection{Почему install \textrightarrow paid conversion = 6\%}

Это один из самых важных assumptions в нашей модели.

\begin{greenbox}[\textbf{Ключевой assumption}]
\[
\text{Install \textrightarrow Paid Conversion} = 6\%
\]
\end{greenbox}

То есть из 100 установок:
\[
100 \times 6\% = 6 \text{ платящих пользователей}
\]

Почему мы считаем это уместным:
\begin{itemize}
    \item высокий intent: беременность --- сильный pain-driven use case;
    \item продукт решает тревогу, а не просто развлекает;
    \item у нас strong premium promise: AI monitoring, alerts, doctor reports;
    \item paywall можно строить вокруг реальной ценности, а не вокруг cosmetic features.
\end{itemize}

Но мы понимаем, что именно conversion --- уязвимое место любой early-stage модели. Поэтому в Excel на листе \textbf{Sensitivity} мы отдельно показываем, как при снижении conversion экономика быстро ухудшается.

\subsection{Почему paid share = 55\%, а organic/referral share = 45\%}

\textbf{Это доля платящих пользователей по источнику привлечения}.

\[
55\% \text{ платящих пользователей приходят из paid}
\]
\[
45\% \text{ платящих пользователей приходят из organic/referrals/doctors/creators}
\]

Почему так:
\begin{itemize}
    \item в начале paid нужен для скорости и тестов;
    \item но health-продукт опасно строить только на платной рекламе;
    \item доверие, рекомендации, word-of-mouth и doctor validation здесь реально важны;
    \item без органики blended CAC будет слишком тяжелым.
\end{itemize}

Именно это и показывает наш лист \textbf{GTM\_Plan}: мы специально не делаем вид, что весь рост идет только через ads.

\subsection{Почему average monetized months = 7}

Это число используется для \textbf{payback}, а не для gross revenue monthly-plan users.  
Тут легко запутаться, поэтому важно разделять:

\begin{itemize}
    \item \textbf{6 месяцев} --- это сколько в среднем платит monthly subscriber;
    \item \textbf{7 месяцев} --- это effective monetized window, по которому мы размазываем LTV для расчета payback.
\end{itemize}

Почему 7 месяцев:
\begin{itemize}
    \item пользователь не обязательно приходит в первый день беременности;
    \item кто-то подключается позже;
    \item часть ценности остается и на раннем postpartum этапе;
    \item мы берем не максимальный срок, а рабочий средний горизонт монетизации.
\end{itemize}

То есть это не ``биологическая длина беременности'', а \textbf{модельный горизонт возврата CAC}.

\section{Как считается revenue и LTV}

В Excel это лежит на листе \textbf{Unit\_Economics}. Ниже та же логика в явном виде.

\subsection{Weighted gross revenue per paid user}

Один платящий пользователь в нашей модели не всегда одинаковый:
\begin{itemize}
    \item 80\% берут annual;
    \item 20\% берут monthly.
\end{itemize}

Поэтому gross revenue считаем как взвешенное среднее:

\[
\text{Revenue} =
(\text{Annual Share} \times \text{Annual Price})
+
(\text{Monthly Share} \times \text{Monthly Price} \times \text{Months on Monthly})
\]

Подставляем:
\[
\text{Revenue} =
0.8 \times 159 + 0.2 \times 15 \times 6
\]

\[
= 127.2 + 18 = \textbf{145.2}
\]

Смысл числа \$145.2:  
это \textbf{средний gross revenue на одного платящего пользователя}, если смотреть на смешанную базу annual + monthly.

\subsection{Net revenue after fee}

Из gross revenue мы вычитаем platform fee:

\[
\text{Net Revenue} = 145.2 \times (1 - 0.15)
\]

\[
= 145.2 \times 0.85 = \textbf{123.42}
\]

\subsection{Contribution LTV}

Теперь вычитаем variable cost per paid user:

\[
\text{LTV} = 123.42 - 9 = \textbf{114.42}
\]

Важно: здесь мы считаем именно \textbf{contribution LTV}, а не просто revenue LTV.  
То есть мы уже убрали переменные затраты и считаем более честную экономику.

\begin{orangebox}[\textbf{Почему это консервативная модель}]
Мы сознательно \textbf{не включаем repeat pregnancies} в базовый LTV. Да, в реальной жизни часть пользователей может вернуться в следующую беременность. Но для competition-логики нам выгоднее показывать более строгую модель: один жизненный цикл пользователя без учета повторных беременностей и возврата к приложению.
\end{orangebox}

\section{Как считается CAC}

\subsection{Paid CAC}

Если установка стоит \$3.8, а в оплату конвертируется 6\% установок, тогда стоимость одного paid user через paid traffic:

\[
CAC_{paid} = \frac{3.8}{0.06} = \textbf{63.33}
\]

Это означает:  
чтобы получить одного платящего пользователя только из paid, нам нужно в среднем потратить \$63.33.

\subsection{Почему blended CAC ниже}

Но реальный бизнес не должен жить только на paid.  
Поэтому мы используем blended model:

\begin{itemize}
    \item 55\% платящих пользователей --- paid;
    \item 45\% --- organic/referral/doctors/creators.
\end{itemize}

Тогда blended CAC считается так:

\[
CAC_{blended} = CAC_{paid} \times \text{Paid Share}
\]

\[
= 63.33 \times 0.55 = \textbf{34.83}
\]

\bluequote{\textbf{Экономика Preggo работает не потому, что paid magically super cheap, а потому что paid усиливается органикой и доверием.}}

\section{Главные итоговые метрики}

\subsection{LTV / CAC}

\[
\frac{114.42}{34.83} = \textbf{3.28x}
\]

Пояснение коэффицентов:
\begin{itemize}
    \item если метрика ниже 1x --- бизнес теряет деньги;
    \item около 2x --- уже что-то есть, но масштабировать опасно;
    \item \textbf{3x+} --- это уже выглядит как отличная база для роста.
\end{itemize}

Наше значение:
\[
\textbf{3.28x}
\]

Это значит, что модель выглядит \textbf{масштабируемой}, если мы удерживаем conversion и annual mix.

\subsection{Payback period}

Сначала считаем monthly contribution:

\[
\text{Monthly Contribution} = \frac{114.42}{7} = \textbf{16.35}
\]

Потом payback:

\[
\text{Payback} = \frac{34.83}{16.35} = \textbf{2.13 \text{ мес}}
\]

Примерно за 2.1 месяца мы отбиваем стоимость привлечения пользователя.

Это сильный результат, потому что:
\begin{itemize}
    \item капитал возвращается быстро;
    \item деньги можно реинвестировать;
    \item модель не требует сверхдолгого ожидания окупаемости.
\end{itemize}

\begin{greenbox}[\textbf{Главный вывод по метрикам}]
\[
\boxed{\text{LTV = \$114.42}}
\qquad
\boxed{\text{CAC = \$34.83}}
\qquad
\boxed{\text{LTV/CAC = 3.28x}}
\qquad
\boxed{\text{Payback = 2.13 мес}}
\]
\end{greenbox}

\section{Что в модели самое чувствительное}

Вот тут как раз нужен лист \textbf{Sensitivity}. Он показывает: не все assumptions одинаково важны.

\begin{redbox}[\textbf{Самые критичные переменные}]
\begin{enumerate}
    \item \textbf{Annual mix}  
    если доля annual падает, revenue и LTV сразу проседают;

    \item \textbf{Install \textrightarrow Paid conversion}  
    если conversion падает ниже base-case, paid CAC резко ухудшается;

    \item \textbf{Paid share}  
    если органика не заводится, blended CAC становится тяжелее;

    \item \textbf{CPI}  
    если аукцион дороже, то paid growth становится менее комфортным.
\end{enumerate}
\end{redbox}

Именно поэтому в нашей модели есть три сценария:
\begin{itemize}
    \item base;
    \item conservative;
    \item optimistic.
\end{itemize}

Мы понимаем, где модель может просесть, и заранее это проверяем.

\section{Как это связано с GTM}

Лист \textbf{GTM\_Plan} отвечает на вопрос из brief'а:  
\textbf{откуда именно будут приходить пользователи и почему эти каналы подходят нам.}

\begin{softbox}[\textbf{Логика каналов роста}]
\begin{itemize}
    \item \textbf{Paid social ads} --- быстрый тест messaging, креативов и paywall;
    \item \textbf{Organic referrals} --- снижают blended CAC;
    \item \textbf{Doctors / clinics} --- добавляют trust;
    \item \textbf{Creators / mom communities} --- дают social proof;
    \item \textbf{ASO / search content} --- ловят already-existing intent.
\end{itemize}
\end{softbox}

То есть наша модель пытается выиграть за счет сочетания:
\[
\text{value} + \text{trust} + \text{organic reinforcement}
\]


\section{Competitor analysis и почему Preggo может выиграть}

Для Preggo эта секция, как и для любого другого стартапа важна, потому что мы не выходим в пустой рынок --- мы выходим в большую категорию pregnancy / women’s health apps, где уже есть сильные игроки, но большинство из них продают в первую очередь \textbf{контент, трекинг, weekly updates, community и визуализацию развития ребенка}, а не систему раннего замечания рисков и doctor-ready monitoring.

\newpage

\begin{softbox}[\textbf{Ключевая логика нашего позиционирования}]
Мы не пытаемся быть ``еще одним приложением про беременность''.  
Мы сознательно занимаем более узкую и более дорогую позицию:

\begin{itemize}
    \item не \textbf{content-first}, а \textbf{safety-first};
    \item не \textbf{``what week am I on?''}, а \textbf{``what is happening with my body right now?''};
    \item не \textbf{развлечение и вовлечение}, а \textbf{monitoring, interpretation, alerts и doctor-facing reports}.
\end{itemize}
\end{softbox}

\subsection{Direct competitors}

\begin{center}
\renewcommand{\arraystretch}{1.25}
\begin{tabularx}{\textwidth}{|p{2.4cm}|p{2.5cm}|p{3.3cm}|p{3.4cm}|X|}
\hline
\rowcolor{lightblue}
\textbf{Игрок} & \textbf{Публичный pricing anchor} & \textbf{Что продает пользователю} & \textbf{Сильные стороны} & \textbf{Где есть окно для Preggo} \\
\hline
\textbf{Flo} & Premium subscription; публично встречаются monthly/yearly offers, в том числе annual around \$59.99 и monthly around \$11.49 в US listing / market intelligence & cycle + fertility + pregnancy tracking, symptom logging, health content, privacy layer & очень сильный бренд, massive scale, привычка пользователя, сильный контентный движок, privacy positioning & Flo --- широкий women’s health product, а не pregnancy-specific safety system. Он силен как broad companion, но не строится вокруг lab scanning, doctor reports и continuous pregnancy risk layer. Также в последнее время юзеры жаловались на проблемы с сохранностью данных и на их продажу, что подпортило доверие компании \\
\hline
\textbf{Pregnancy+} & freemium / premium / in-app purchases & weekly pregnancy content, expert articles, baby development visuals, interactive 3D models & очень сильный emotional UX, визуализация ребенка, понятный mass-market продукт, huge download base & продукт в основном про guidance, content и visualization. Это очень хороший emotional tracker, но не полноценная decision-support система для риска и интерпретации данных \\
\hline
\textbf{Glow Nurture} & Glow Premium publicly positioned around \$59.99/year & pregnancy tracking, personalized tips, reminders, community, AI-style insights / comparisons & экосистема Glow, community, strong habit loop, familiar journaling/logging model & даже при AI wording core product все еще ближе к tracking + support + community, чем к medically structured monitoring и doctor-facing output \\
\hline
\textbf{What to Expect} & в публичном positioning упор больше на trusted content/community, чем на premium medical layer & weekly updates, due date tracking, articles, baby-size metaphors, community/forums & сильный media brand, доверие, SEO/content moat, огромный top-of-funnel & отличный informational brand, но не deep monitoring product. Это скорее education + community layer, чем safety infrastructure \\
\hline
\textbf{Preggo} & \textbf{\$159 annual / \$15 monthly} & continuous monitoring, AI risk logic, lab/document scanner, alerts, doctor report & более сильная value density на одну беременность; clearer premium logic; closer to actionable healthcare workflow & наш риск --- выше цена и уже аудитория. Но именно за счет этого у нас выше ARPU и право на более сильную unit economics модель \\
\hline
\end{tabularx}
\end{center}

\subsection{Связь с pricing}

Наш pricing выше, чем у типичных контентных pregnancy apps, и это \textbf{осознанно}. Мы не обязаны стоить как mass-market tracker, если даем продукт другого класса.

\begin{orangebox}[\textbf{Почему premium pricing уместен}]
Большинство существующих игроков монетизируют:
\begin{itemize}
    \item контент;
    \item эмоциональную вовлеченность;
    \item трекинг недели беременности;
    \item community / education.
\end{itemize}

Preggo монетизирует более дорогую ценность:
\begin{itemize}
    \item снижение тревоги между визитами к врачу;
    \item интерпретацию симптомов и биоданных;
    \item раннее замечание отклонений;
    \item удобную упаковку данных для визита к врачу.
\end{itemize}

То есть \$159 annual --- это не ``дорогой pregnancy content app'', а \textbf{дешевый по сравнению с ценой неопределенности, лишних визитов, missed signals и когнитивной перегрузки во время беременности}.
\end{orangebox}

\subsection{Marketing competitor analysis}

Даже маркетинг конкурентов показывает, что категория в основном продается через более мягкие углы:

\begin{itemize}
    \item \textbf{Flo} продвигает privacy, symptom tracking, personalized insights и broad women’s health utility;
    \item \textbf{Pregnancy+} делает акцент на beautiful 3D baby development, weekly guidance, expert content и emotional support;
    \item \textbf{Glow Nurture} продвигает personalized pregnancy tracking, community и ``Is this normal?''-style reassurance;
    \item \textbf{What to Expect} исторически очень силен через trusted content, week-by-week guidance и community brand.
\end{itemize}

\bluequote{\textbf{Общий паттерн рынка: большинство конкурентов продают информацию, привычку и эмоциональную поддержку. Preggo должен продавать спокойствие через actionable monitoring.}}

Из этого для нас следует конкретная marketing-логика:

\begin{greenbox}[\textbf{Какие message angles разумно использовать Preggo}]
\begin{enumerate}
    \item \textbf{``Know earlier, worry less''}  
    не просто ``следи за беременностью'', а ``раньше замечай сигналы риска'';

    \item \textbf{``From scattered symptoms to one clear picture''}  
    не ручной хаос и заметки, а единая картина состояния;

    \item \textbf{``Doctor-ready, not just mom-content''}  
    не только статьи и baby-size metaphors, а данные, которые можно показать врачу;

    \item \textbf{``Your pregnancy, interpreted''}  
    не просто logging, а интерпретация и приоритизация того, на что реально стоит обратить внимание.
\end{enumerate}
\end{greenbox}

\subsection{Почему это усиливает нашу экономику}

Competitor analysis важен не сам по себе, а потому что он объясняет, \textbf{почему наш ARPU и annual-first model вообще имеют право на существование}.

\begin{itemize}
    \item если бы мы были еще одним weekly-content tracker, то \$159 annual выглядел бы слишком агрессивно;
    \item но если продукт действительно решает pain of uncertainty, improves monitoring и экономит ментальную нагрузку, higher willingness-to-pay становится логичной;
    \item именно поэтому наша unit economics строится не на mass-market entertainment logic, а на \textbf{higher-value, narrower-segment monetization}.
\end{itemize}

\begin{tcolorbox}[
    enhanced,
    breakable,
    colback=brandblush,
    colframe=brandcoraldeep,
    boxrule=0.85pt,
    arc=4pt,
    left=10pt,right=10pt,top=8pt,bottom=8pt
]
\textbf{Финальный вывод по конкурентам.}
Рынок не пустой, но он занят в основном продуктами, которые отлично объясняют беременность, сопровождают ее и вовлекают пользователя через контент, комьюнити и визуализацию. Preggo идет не в лобовую войну за ``лучший контент'', а в более защищенную позицию: \textbf{pregnancy intelligence + monitoring + risk interpretation + doctor-facing output}. Именно это и делает наш premium pricing и наши unit economics более защищенными.
\end{tcolorbox}

\vspace{4pt}
{\footnotesize \textcolor{mutedtext}{Источники для competitor block: Flo official product/help pages, Flo market pricing listings, Philips Pregnancy+ official pages, Glow official site and Nurture pages, What to Expect official app listings, Meta Ads Library public page listings. Цены и офферы могут отличаться по стране, платформе и промо-периоду, поэтому в модели они используются как публичные рыночные anchors, а не как immutable constants.}}

\section{Почему мы считаем модель реалистичной}

Мы специально не делали модель сказочной. Наоборот:

\begin{itemize}
    \item не включали repeat pregnancies в base LTV;
    \item не ставили слишком низкий CPI;
    \item не ставили absurdly high conversion;
    \item оставили variable costs;
    \item показали sensitivity, а не одну ``идеальную'' картинку.
\end{itemize}

Мы не говорим, что все будет идеально. Мы говорим, что даже при достаточно собранной и трезвой модели Preggo уже выглядит как продукт с отличной unit economics.

\section{Финальный вывод}

Экономика Preggo строится на трех вещах:

\begin{enumerate}
    \item \textbf{Annual-first pricing}  
    потому что use case длинный и ценность высокая;

    \item \textbf{Сильный premium promise}  
    потому что мы продаем не просто трекинг, а спокойствие, интерпретацию и раннее замечание рисков;

    \item \textbf{Blended acquisition}  
    потому что в health-категории нельзя жить только на paid ads.
\end{enumerate}

\vspace{3pt}

\begin{tcolorbox}[
    enhanced,
    breakable,
    colback=brandblush,
    colframe=brandcoraldeep,
    boxrule=0.9pt,
    arc=5pt,
    left=12pt,right=12pt,top=10pt,bottom=10pt,
    borderline west={3pt}{0pt}{brandcoraldeep}
]
{\color{brandcoraldeep}\Large\bfseries Финальная позиция}\\[6pt]
{\color{darktext}
Наша логика простая: Preggo --- это не просто красивая идея, а продукт, у которого сходится экономика. Мы понимаем, откуда берутся пользователи, как они монетизируются, какие assumptions самые важные и где находятся главные риски. А значит, это уже не просто концепт --- это бизнес, который можно защищать, тестировать и масштабировать.
}
\end{tcolorbox}

\vfill
\begin{center}
\textcolor{mutedtext}{Kazaryan Leon, Adlet Dautbek}
\end{center}

\end{document}